\documentclass[12pt]{article}
\usepackage[utf8]{inputenc}
\usepackage[margin=1in]{geometry}
\usepackage{setspace}
\usepackage{cite}
\usepackage{url}
\usepackage{hyperref}

\doublespacing

\title{Literature Review: Software Metrics in Continuous Integration
and Continuous Delivery Pipelines}
\author{CSE 7319 - Software Metrics and Quality}
\date{\today}

\begin{document}

\maketitle

\section{Introduction}

This literature review examines recent research on software metrics
applied to Continuous Integration and Continuous Delivery (CI/CD)
pipelines, using the fundamental distinction between internal and
external metrics as an organizing framework. Internal metrics—which
depend only on program artifacts and are available earlier in
development with finer granularity—provide controllable indicators of
code quality, complexity, and test adequacy. External metrics—which
reflect customer/user perspectives and system "-ilities" such as
reliability, performance, and delivery efficiency—represent the
ultimate quality goals but are measured later and are harder to
control directly. The fundamental relationship is predictive: internal
metrics serve as independent variables (leading indicators) that help
predict and control external metrics (dependent variables or targets).

CI/CD environments provide an ideal context for applying the
Goal-Question-Metric (GQM) framework, where goals are defined in terms
of external quality attributes, questions probe whether goals are
being met, and metrics (both internal and external) provide
quantitative answers. The selected papers demonstrate how modern
DevOps practices enable automated collection, analysis, and
visualization of both metric types, supporting the quality improvement
paradigm through continuous feedback loops.

The papers reviewed were published between 2016 and 2025, with
emphasis on recent work from the past five years. All sources are
peer-reviewed publications from IEEE, ACM, and other reputable academic venues.

\section{Internal Metrics}

Internal metrics depend solely on program artifacts and can be
measured through static or dynamic analysis. As independent variables
in the metrics relationship model, they serve as leading indicators
and control mechanisms for achieving external quality goals. Classic
categories include size metrics (LOC, function points), complexity
metrics (cyclomatic complexity, Halstead measures, data/control flow),
and structure metrics (coupling, cohesion). The papers reviewed here
focus on how internal metrics are collected and applied within CI/CD
pipelines, particularly static analysis quality checks and test
coverage measurements.

\subsection{Static Code Analysis Metrics}

\subsubsection{How Open Source Projects Use Static Code Analysis
Tools in Continuous Integration Pipelines}

\textbf{Citation:} \cite{zampetti2017static}

\textbf{Summary:} This empirical study investigates how open-source
projects integrate static code analysis tools into their continuous
integration workflows. The researchers analyzed real-world CI
configurations to understand what types of code quality issues cause
builds to fail versus generate warnings, and how developers respond
to these signals. The study examined projects using tools like
Checkstyle, PMD, and FindBugs integrated into CI systems like Travis
CI and Jenkins. A key finding was that the vast majority of build
failures triggered by static analysis concerned coding standard
violations and style issues rather than actual bug detection. The
researchers also found that when static analysis generates warnings
or failures, developers typically address the issues directly by
fixing the code rather than suppressing warnings or adjusting tool
configurations. This suggests that when properly integrated, static
analysis in CI pipelines can effectively enforce code quality standards.

\textbf{Analysis:} This paper demonstrates the practical application
of internal metrics (specifically presentation and structure
complexity indicators from static analysis tools) as control
mechanisms within CI/CD pipelines. The finding that coding standards
violations dominate build failures reflects teams using these internal
metrics to enforce maintainability—one of the key external quality
attributes. This illustrates the predictive relationship: by
controlling internal metrics (code style, structural consistency),
teams aim to improve external metrics (maintainability effort,
long-term productivity). The paper validates that automated quality
gates can effectively enforce internal metric thresholds, embodying
the measurement-to-action workflow emphasized in the quality
improvement paradigm. The developers' tendency to fix issues rather
than suppress warnings demonstrates that when internal metrics are
properly integrated into workflows, they influence behavior and drive
quality improvements. From a GQM perspective, these teams have defined
maintainability goals, posed questions about code consistency, and
selected metrics (style violations, complexity checks) that automated
tools measure continuously.

\subsection{Test Coverage Metrics}

\subsubsection{Open Code Coverage Framework: A Consistent and
Flexible Framework for Measuring Test Coverage}

\textbf{Citation:} \cite{sakamoto2010open}

\textbf{Summary:} This paper addresses a fundamental challenge in
measuring test coverage across different programming languages and
coverage criteria: the high cost and inconsistency of coverage
measurement tools. The researchers observed that while test coverage
is a critical quality indicator, each language typically requires
separate tools with inconsistent reporting formats and capabilities.
They developed the Open Code Coverage Framework (OCCF), which
abstracts common concepts across programming languages using abstract
syntax trees (ASTs). By identifying structural commonalities at the
AST level, the framework can support multiple languages with
significantly less implementation effort. The authors demonstrated
that their approach reduced implementation requirements by
approximately 90\% and cut development time for new coverage criteria
by 80\% or more. The framework supports various coverage types
including statement, branch, and path coverage while remaining
extensible for custom criteria.

\textbf{Analysis:} This paper addresses a fundamental challenge in
measurement theory: ensuring consistent and properly defined metrics
across different contexts. Test coverage is an internal metric (it
depends only on code and test artifacts), and serves as a leading
indicator for external quality attributes like reliability and defect
density. However, the predictive relationship only holds when coverage
is measured consistently using well-defined criteria. The framework's
abstraction approach—using ASTs to identify structural commonalities—
enables uniform measurement of statement, branch, and path coverage
across languages. This consistency is essential for establishing valid
predictive models: if coverage is measured differently across
components, correlations with external quality metrics become
unreliable. The significant reduction in implementation effort (90\%)
demonstrates practical value, making comprehensive internal metrics
collection feasible in heterogeneous environments. From an evaluation
perspective, this framework is properly defined (clear criteria),
properly used (consistent application), and produces useful results
(validated coverage measurements), meeting the key criteria for
effective software metrics.

\section{External Metrics}

External metrics reflect quality from the customer and user
perspective, measuring the "-ilities" (reliability, usability,
performance, security, maintainability) that determine system value.
As dependent variables in the metrics relationship model, these are
the ultimate targets—but they are measured late in the development
cycle and are difficult to control directly. Instead, they must be
predicted and influenced through careful management of internal
metrics. In CI/CD contexts, external metrics span both traditional
quality attributes (reliability, performance) and process efficiency
measures (deployment frequency, lead time, change failure rate). The
papers reviewed demonstrate approaches to measuring and improving
these outcome-oriented metrics through automated data collection in
DevOps pipelines.

\subsection{DevOps Research and Assessment (DORA) Metrics}

\subsubsection{A Framework for Automating the Measurement of DORA Metrics}

\textbf{Citation:} \cite{wilkes2023framework}

\textbf{Summary:} The DORA metrics (Deployment Frequency, Lead Time
for Change, Time to Restore Service, and Change Failure Rate) have
become industry-standard indicators of software delivery performance,
but many organizations struggle to measure them consistently or
resort to manual data collection. This paper presents a framework
that automatically extracts and calculates DORA metrics from project
repositories without requiring specific CI/CD tools or cloud
platforms. The researchers validated their framework by analyzing 304
popular open-source projects on GitHub, generating time-series data
that revealed meaningful insights into project performance trends. A
key contribution is the project-agnostic design: the framework can
operate across diverse development lifecycles and deployment
approaches. The analysis showed that significant changes in DORA
metrics correlated with major historical events in projects,
validating the metrics' ability to detect meaningful shifts in
development performance. The authors recommend augmenting automated
measurements with additional contextual data like bug criticality
when available.

\textbf{Analysis:} The DORA metrics represent external quality
attributes in the DevOps context, mapping to classical metrics
categories: deployment frequency and lead time reflect productivity
and schedule performance, while change failure rate and time to
restore service measure reliability. These align with the Putnam/Myers
framework's external metrics (time, quality, productivity). The
framework's project-agnostic design is significant because it enables
consistent measurement across diverse environments, satisfying the
"properly defined" and "properly used" evaluation criteria. Critically,
this paper demonstrates that external metrics can be extracted
automatically from development artifacts (commits, deployments,
incidents), rather than requiring manual surveys or estimation—
addressing the traditional challenge that external metrics are "hard
to control" and "measured late." The time-series approach is essential
for quality improvement: tracking metric trends over time enables
teams to assess whether process changes (such as adopting new internal
metric thresholds or tools) actually improve external outcomes. The
correlation with historical project events validates that these
measurements capture meaningful quality signals rather than arbitrary
numbers, meeting the "useful results" criterion.

\subsection{CI/CD Metrics Visualization}

\subsubsection{A Roadmap for Using Continuous Integration Environments}

\textbf{Citation:} \cite{yu2024roadmap}

\textbf{Summary:} This recent study investigates how development
teams can leverage CI environments to generate and visualize quality
metrics. Through interviews with 22 participants across five
industrial projects, the researchers identified that 86\% of
participants found CI-generated data useful for quality attribute
evaluation, and 73\% valued centralized dashboards for visualizing
this data. The study cataloged 14 base metrics and 10 derived metrics
spanning seven quality attributes including maintainability,
security, performance, and reliability. Performance metrics showed
the highest adoption at 82\% among participants. The paper provides
concrete examples of metrics collection pipelines using Prometheus
for data gathering and Grafana for visualization. The researchers
also propose a three-level maturity model (elementary, intermediate,
advanced) for CI environment evolution, finding that 60\% of studied
projects operated at intermediate maturity levels. Challenges
identified include dashboard information overload and insufficient
understanding of metrics meanings.

\textbf{Analysis:} This paper directly addresses the metrics usage
workflow: after measurement activities collect data, analysis and
visualization are essential to make results actionable. The study
validates the GQM framework in practice: teams use CI-generated data
to evaluate quality attributes (goals), with 86\% finding it useful
for this purpose. The cataloging of 14 base and 10 derived metrics
spanning seven quality attributes demonstrates the breadth of both
internal (maintainability, security) and external (performance,
reliability) metrics available in modern CI/CD systems. Performance
metrics showing the highest adoption (82\%) reflects their direct
connection to external user experience, while the centralized
dashboard preference (73\%) confirms that integrated visualization of
multiple metrics is more valuable than isolated measurements. The
identified challenges—information overload and insufficient
understanding—highlight a critical gap between data collection and
effective use. This suggests the need for "Experience Factory"
approaches where organizations learn which metrics matter and how to
interpret them. The three-level maturity model recognizes that metrics
programs evolve: elementary levels focus on basic measurement,
intermediate levels on integration and automation, and advanced levels
on predictive analytics and optimization—mirroring the progression
from simple measurement toward decision support systems.

\subsection{Requirements Traceability}

\subsubsection{Continuous Integration and Delivery Traceability in
Industry: Needs and Practices}

\textbf{Citation:} \cite{stahl2016continuous}

\textbf{Summary:} This paper addresses the challenge of maintaining
requirements traceability in fast-paced CI/CD environments where
continuous changes make traditional traceability approaches
impractical. The authors note that while traceability has long been
recognized as important for legal, certification, and quality
assurance purposes, empirical studies of traceability practices in
industry—especially regarding tooling and infrastructure—are limited.
The paper documents the Eiffel framework, developed in industry to
provide real-time traceability throughout CI/CD pipelines. Unlike
traditional traceability matrices that are maintained manually or
generated periodically, the Eiffel approach embeds traceability into
the continuous integration infrastructure itself, automatically
capturing relationships between requirements, commits, builds, tests,
and deployments as they occur. This enables immediate visibility into
which requirements are affected by any change and what testing has
been performed.

\textbf{Analysis:} Requirements traceability represents a hybrid
metric that bridges internal and external quality measurement. It
connects external requirements (customer needs, specified
functionality) to internal artifacts (commits, builds, tests),
enabling impact analysis and verification. The Eiffel framework
demonstrates automated measurement embedded directly into development
workflows—addressing the classic challenge that manual metrics
collection is impractical in fast-paced development. This automation
enables real-time tracking rather than periodic snapshots, making
traceability a useful leading indicator: when a requirement changes,
teams immediately know which code and tests are affected. The
framework supports different stakeholder perspectives (developer vs.
manager views), recognizing that effective metrics must answer
specific questions for specific goals—the essence of GQM. From a
measurement theory perspective, this automated approach produces more
reliable data than manual maintenance because it captures actual
relationships (commits linked to requirements) rather than documented
intentions. The integration of traceability with CI/CD metrics (build
status, test results, deployment records) enables holistic quality
assessment: teams can track not just whether code builds successfully
(internal metric) but whether specific requirements are fully
implemented and tested (external validation).

\subsection{Hybrid Development Models}

\subsubsection{Mathematical Model of the Software Development Process
with Hybrid Management Elements}

\textbf{Citation:} \cite{semenov2025mathematical}

\textbf{Summary:} This very recent paper develops a probabilistic
model for hybrid software development lifecycles that combines
traditional project management with agile and DevOps practices. The
model integrates data from multiple sources including requirements
traceability systems, static code analysis tools, and CI/CD pipeline
telemetry. Using GERT (Graphical Evaluation and Review Technique)
network semantics and fuzzy uncertainty treatment, the model predicts
completion timelines and rework requirements. A particularly novel
aspect is the incorporation of AI-assisted development into the
model, allowing comparison of development scenarios with and without
AI tooling. The validation demonstrated approximately a threefold
reduction in expected rework cycles when AI assistance was enabled,
along with earlier defect detection. The framework yields
interpretable metrics for tuning quality gates and automation levels
in Agile/DevOps settings.

\textbf{Analysis:} This paper exemplifies the predictive relationship
between internal and external metrics through formal modeling. By
integrating data from requirements traceability (external
specifications), static code analysis (internal quality metrics), and
CI/CD telemetry (both internal build metrics and external delivery
performance), the model establishes quantitative relationships between
leading indicators and outcome variables. The probabilistic approach
using GERT networks addresses the fundamental challenge in software
metrics: external quality attributes like schedule and rework are not
uniquely determined by internal metrics, but can be predicted with
confidence intervals. The threefold reduction in rework cycles with AI
assistance demonstrates how changes in development processes alter the
relationships between metrics—requiring recalibration of predictive
models as practices evolve. This aligns with the Quality Improvement
Paradigm: measure baseline performance, introduce changes (AI tools),
quantify impacts, and update understanding. The focus on tuning
quality gates shows metrics used not just for assessment but for
active control—deciding which internal metric thresholds to enforce
based on their predicted impact on external outcomes. This represents
the evolution from basic measurement toward integrated decision
support, where multiple metric types combine to provide actionable
guidance rather than isolated data points.

\section{Conclusion}

The reviewed papers demonstrate how fundamental software metrics
principles apply in modern CI/CD environments. Several key themes
connect the research to metrics theory:

\textbf{Internal and External Metrics Integration:} The papers
illustrate the predictive relationship between internal metrics
(static analysis violations, test coverage, code complexity) and
external quality outcomes (maintainability, reliability, delivery
performance). CI/CD pipelines provide infrastructure for measuring
both metric types automatically and establishing their correlations,
as demonstrated in the DORA framework and Semenov's probabilistic
model.

\textbf{GQM Framework in Practice:} The Yu et al. study validates the
Goal-Question-Metric approach: teams define quality attribute goals
(performance, reliability, security), pose questions about current
state, and select specific metrics to answer them. The challenge of
metrics comprehension highlights that effective GQM requires not just
data collection but clear understanding of what metrics mean and how
they relate to goals.

\textbf{Automation Enables Measurement Theory:} Manual collection
violates the principle that metrics must be properly defined and
consistently applied. The automated frameworks reviewed (DORA, Eiffel,
OCCF) ensure uniform measurement across diverse contexts, enabling
reliable predictive models. This addresses the classical challenge
that inconsistent measurement produces unreliable correlations between
internal and external metrics.

\textbf{From Measurement to Control:} The papers demonstrate metrics
usage workflow: measurement activities (automated collection) →
analysis (dashboards, time-series visualization, predictive models) →
action (quality gates, process improvements) → validation (experience
factory learning). This progression shows metrics evolving from simple
descriptive statistics toward active control mechanisms that guide
development decisions.

\textbf{Holistic Quality Assessment:} Integrating multiple metric
types—internal code quality, test adequacy, requirements traceability,
and external delivery performance—provides richer insights than
isolated measurements. The Semenov model exemplifies this by
synthesizing requirements, static analysis, and CI/CD data into
unified predictions.

These papers confirm that CI/CD environments realize the vision of
comprehensive, automated metrics programs. They demonstrate that when
properly implemented, metrics frameworks support continuous quality
improvement by making both internal leading indicators and external
outcome measures visible, actionable, and connected through validated
predictive relationships.

\bibliographystyle{IEEEtran}
\bibliography{references}

\end{document}
